% IACR Transactions CLASS DOCUMENTATION
% Written by Gaetan Leurent gaetan.leurent@inria.fr (2016-2018)
%
% To the extent possible under law, the author(s) have dedicated all
% copyright and related and neighboring rights to this software to the
% public domain worldwide. This software is distributed without any
% warranty.
%
% You should have received a copy of the CC0 Public Domain Dedication
% along with this software. If not, see
% <http://creativecommons.org/publicdomain/zero/1.0/>.

\documentclass{iacrtrans}
\usepackage[utf8]{inputenc}

\usepackage{booktabs}

\author{Gaëtan Leurent\inst{1} \and Friedrich Wiemer\inst{2}}
\institute{Inria, Paris, France, \email{gaetan.leurent@inria.fr} \and
           Ruhr-Universität Bochum, Bochum, Germany, \email{friedrich.wiemer@rub.de}}
\title[\texttt{iacrtans} class documentation]{\publname}
\subtitle{\LaTeX{} Class Documentation (v. 0.92)}

\begin{document}

\maketitle

% use optional argument because the \LaTeX command breaks the PDF keywords
\keywords[\publname, ToSC, TCHES, LaTeX]{\publname{} \and ToSC
  \and TCHES \and \LaTeX}

\begin{abstract}
  This document is a quick introduction to the \LaTeX{} class for the
  \publname{}.
\end{abstract}

\tableofcontents{}

\section*{Introduction}

The \texttt{iacrtans} \LaTeX{} class is used by the new ``IACR
Transactions'' journals ToSC (IACR Transactions on Symmetric Cryptology)
and TCHES (IACR Transactions on Cryptographic Hardware and Embedded
Systems).  The class is based on standard \LaTeX{} classes and packages
(mainly the \texttt{article} class with \texttt{amsmath}), and should be
similar in use to the \texttt{llncs} class used for Springer's
proceedings.

The class is still in development and feedback and comments are welcome.
The latest version can be found on the Github page of the project:
\url{https://github.com/Cryptosaurus/iacrtrans}, feel free to open
tickets for issues or to submit pull requests.

\section{\textcolor{red!70!black}{FAQ}: Frequently Asked Questions}

\subsection{Converting \texttt{llncs} Papers to \texttt{iacrtrans} for Submission}

If you have a paper typeset with the \texttt{llncs} class, conversion
should be relatively easy.  The following steps should be sufficient
in most cases (for the submission version):
\begin{enumerate}
\item Replace \verb+\documentclass{llncs}+ with\\
    \verb+\documentclass[journal=XXX,submission,spthm]{iacrtrans}+,
    where \verb+XXX+ is either \verb+tosc+, or \verb+tches+;
\item Replace \verb+\bibliographystyle{splncs03}+ with
  \verb+\bibliographystyle{alpha}+;
\item Add a \verb+\keywords{}+ command before the abstract, with
  keywords separated by \verb+\and+;
\item Remove commands that might override the class style, such as
  \verb+\pagestyle{...}+ or \verb+\thispagestyle{...}+,
  change of margins (\emph{e.\,g.}  with the \texttt{geometry} package),
  change of fonts, \ldots
\item See also \autoref{sec:biblio} for information about how to
  typeset the bibliography.
\end{enumerate}

\subsection{Template File}

There is a \texttt{template.tex} file included with this class, to use
as a starting point for writing a paper.  The \LaTeX{} source of this
documentation can also be used as a more advanced example, but please
make sure to remove any unnecessary code.

\subsection{Preparing the Final Version}

Before submitting your final version, please make sure that it compiles
properly with the \texttt{[final]} option (see the next subsection in
case of issues), and check that the author names and affiliation are
correct.  In particular, each affiliation MUST be of the form
$\langle\text{institue}\rangle, \langle\text{city}\rangle,
\langle\text{country}\rangle$.

\subsection{Compilation Issues}

If your document doesn't compile with the \texttt{iacrtrans} class,
you can try adding the \texttt{[nohyperref]} option.  This will
disable some code in the class that is known to be fragile, in
particular the generation of metadata.  If this fixes your compilation
problems, please try to define a clean version of \texttt{author}, \texttt{institute},
\texttt{title} and/or \texttt{keywords} as optional arguments to
these commands (in particular, remove \LaTeX{} macros, and write
everything on a single line), and see if you document now compiles
without the \texttt{[nohyperref]} option.

\subsection{Line Number Issues}

In submission mode, the class adds line numbers to help reviewers refer
to specific elements.  Unfortunately, the code that does this is rather
fragile, and some constructions break the line numbering.  In
particular, old-style equations with \verb+$$...$$+ are known to cause
issues, please use \verb+\[...\]+ instead\footnote{For more background
  on \texttt{\$\$...\$\$} and \texttt{\textbackslash[...\textbackslash]},
  see \url{https://tex.stackexchange.com/questions/503/why-is-preferable-to}}.

If you find other cases where line numbering is broken, please open a
ticket on github, and we will try to find a solution when possible.
As a quick workaround, you can wrap the offending environment with
\verb+\begin{linenomath}+ / \verb+\end{linenomath}+ (this will do
nothing in preprint or final mode when line numbers are deactivated).

\section{Main Commands}

\subsection{Title Page}

The following commands are used to input informations for the title page.

\paragraph{\texttt{\textbackslash title}} to define the title.

A shorter running title can be given as optional argument.

\paragraph{\texttt{\textbackslash subtitle}} to give an optional
subtitle.

\paragraph{\texttt{\textbackslash author}} to define the author list.

Author names must be delimited by \verb+\and+ macros.  If there is one
different affiliation for each author, authors and affiliations will be
numbered automatically.  Otherwise, each author name must be followed by
\verb+\inst{...}+ with the corresponding affiliation(s).

A shorter list of authors for the running head can be given as
optional argument.

\paragraph{\texttt{\textbackslash institute}} to give author's affiliation(s).

If there are several affiliations, they must be separated by
\verb+\and+ macros, and will be numbered automatically.

\paragraph{\texttt{\textbackslash keywords}} to give a list of
keywords.

Individual keywords should be separated by the \verb+\and+ macro.

If there are fragile commands in the keywords, use the optional argument
to give a text-only version of the keywords; this will be used for the
PDF metadata.

\paragraph{\texttt{\textbackslash email}} should be used inside the
\verb+\institute+ argument to typeset author's email address(es).  An
optional argument can be given for the hyperlink, if different from the
displayed email.  For instance, you can group emails as follows:\\
\verb+\email[alice@foo.com,bob@bob.com]{{alice,bob}@foo.com}+

\paragraph{\texttt{\textbackslash thanks}}
can be used inside the \verb+\title+,
\verb+\author+ or \verb+\institute+ argument to generate a footnote with additional
information, if needed.

\paragraph{\texttt{\textbackslash maketitle}} is used to actually
typeset the title.

\paragraph{The \texttt{abstract} environment} should be used to typeset the abstract.

Note that the keywords should be given before starting the abstract environment.


\subsection{Theorems}

The \texttt{iacrtrans} class uses the \AmS{} packages to typeset
math.  In particular, it loads the \texttt{amsthm} package, and
predefines the following environments:
\begin{center}
  \ttfamily
\begin{tabular}{l@{\hspace{1cm}}l@{\hspace{1cm}}l}
theorem     & definition & remark \\
proposition & example    & note   \\
problem     & exercise   & case   \\
lemma       & property   &        \\
conjecture  & question   &        \\
corollary   & solution   &        \\
claim       &            &        \\
\end{tabular}
\end{center}

Note that the \texttt{proof} environment automatically adds a QED
symbol at the end of the proof (unless you give option
\texttt{[spthm]} to the \texttt{iacrtrans} class).  If the QED symbol
is typeset at a wrong position, you can force its position with
\verb+\qedhere+.

\section{Class Options}
\label{sec:options}

\subsection{Publication Type}

The class supports four publication types, selected with the
following class options:
\begin{description}
\item[\texttt{[final]}] for final papers
\item[\texttt{[preprint]}] for preprints (without copyright info, default)
\item[\texttt{[submission]}] for submissions (anonymous, with line numbers)
\item[\texttt{[draft]}] is similar to preprint, but activates draft
  mode for the underlying \texttt{article} class (which shows overfull hboxes), and other packages
  (\emph{e.\,g.} \texttt{graphicx}, \texttt{hyperref}).
\end{description}
Some types supports further options:
\begin{description}
\item[\texttt{[journal=XXX]}] must be used for submissions and
  final papers; allowed values are \texttt{tosc} and \texttt{tches},
  to set the \texttt{publname} macro accordingly to the right journal
\item[\texttt{[notanonymous]}] can be used in submission mode, when the
  call for paper requires non-anonymous submissions.
\end{description}

\subsection{Other Options}

\paragraph{\texttt{[spthm]}}
provides theorem environments that emulates
\texttt{llncs} class's \texttt{sptheorem}:
\begin{itemize}
\item A \texttt{\textbackslash spnewtheorem} wrapper is provided
  around \AmS{} \texttt{\textbackslash newtheorem}.  Note that the
  styling options are ignored; you should use standard
  \texttt{amsthm} commands for fine control.
\item The \AmS{} \texttt{proof} environment will not automatically add a
  QED symbol at the end of the proof.
\end{itemize}

\paragraph{\texttt{[floatrow]}}
uses the \texttt{floatrow} package to customize floats rather than the
plain \texttt{float} package.  In particular, this allows to typeset
floats side by side as shown in this example:
\begin{verbatim}
\documentclass[floatrow]{iacrtrans}
\usepackage[demo]{graphicx}
\begin{document}

\begin{figure}
  \begin{floatrow}
    \ffigbox{\includegraphics[width=0.4\textwidth]{1.png}}
    {\caption{This is caption 1.}}
    \ffigbox{\includegraphics[width=0.4\textwidth]{2.png}}
    {\caption{This is caption 2.}}
  \end{floatrow}
\end{figure}

\end{document}
\end{verbatim}

The row will be divided equally according to the number of figures, but
you can ask each figure to take its natural space instead with
\verb+\ffigbox[\FBwidth]+.  For more advanced use, see the
\texttt{floatrow} documentation.

\paragraph{\texttt{[xcolor=$\langle\text{\emph{list of options}}\rangle$]}}
passes $\langle\text{\emph{list of options}}\rangle$ to the
\texttt{xcolor} package.  Since \texttt{xcolor} is loaded by the class,
you have to use this mechanism to pass options at load time; for instance
use \texttt{xcolor=svgnames} to load SVG color names.


\paragraph{\texttt{[hyperref=$\langle\text{\emph{list of options}}\rangle$]}}
passes $\langle\text{\emph{list of options}}\rangle$ to the
\texttt{hyperref} package.  Alternatively, you can load
\texttt{hyperref} yourself with the required options and the class will
detect that it already loaded.

\paragraph{\texttt{[nohyperref]}}
disables the automatic loading of
\texttt{hyperref}.  Use this is if your document fails to compile with
\texttt{hyperref} for some reason.

The \texttt{iacrtrans} class automatically loads \texttt{hyperref}
after all other packages.  If you need some packages to be loaded
\emph{after} \texttt{hyperref}, you should load \texttt{hyperref}
explicitly at the correct position, but not use the \texttt{[nohyperref]} option.

% \paragraph{\texttt{[metada]}}
% activates the automatic creation of PDF metadata.  This is
% enabled automatically byin the \texttt{final} version, but it has been disabled by default
% with other because the required sting manipulations are not very robust, and can
% result in weird errors.  If you get errors with this option, you
% should try to define a clean version of \texttt{author},
% \texttt{title} and/or \texttt{keywords} as optionnal arguments to
% these commands.

\paragraph{\texttt{[nolastpage]}}
disables the automatic loading the \texttt{lastpage} package in
\texttt{[final]} mode.  When this option is enabled, the last page
number must be set explicitly with
\texttt{\textbackslash{}setlastpage}.


\section{Typesetting the Bibliography}
\label{sec:biblio}

% Borrowed from btxdoc.tex
\def\BibTeX{{\rm B\kern-.05em{\sc i\kern-.025em b}\kern-.08em
    T\kern-.1667em\lower.7ex\hbox{E}\kern-.125emX}}

Having good bibliographic references is very important for the
visibility of the journal.  Since we don't have a commercial editor,
authors need to make sure themselves that references are standardized
and clean.  We strongly encourage authors to use \BibTeX{} for the
bibliograpy, using bibliographic data
from \url{http://www.dblp.org} or \url{https://cryptobib.di.ens.fr/}.

We are still working on a good solution for the bibliography, and we
expect to have more specific instructions when producing the final
version of the papers, including a dedicated \BibTeX{} style.


\section{Further Instructions}

\paragraph{\LaTeX{} distribution, and worklow.}  \LaTeX{}
distributions are available on a variaty of platforms.  In particular,
we recommand the \href{https://www.tug.org/texlive/}{TeX Live}
distribution, which is updated regularly, include a large number of
packages, and is available on many platforms.
\begin{description}
\item[Linux:] A LaTeX installation is included in most Linux
  distributions.  Alternatively,
  \href{https://www.tug.org/texlive/}{TeX Live} can be installed
  easily without root access.
\item[Windows:] There are also good \LaTeX{} distributions for Windows,
  such as \href{http://www.miktex.org/}{MikTeX} and
  \href{https://www.tug.org/texlive/}{TeX Live}.
\item[MacOSX:] On MacOSX, TeX Live is available inside
  \href{http://www.tug.org/mactex/}{MacTeX}.
\end{description}


We recommand the use of \texttt{pdflatex} because it generally
supports more features than \texttt{latex} and \texttt{dvips}
(\texttt{xelatex} and \texttt{lualatex} are also missing some advanced
features from \texttt{pdflatex}).

\paragraph{Section and paragraph headings.}

Section heading of all levels should be typeset in Camel Case.
%
Paragraph headings (\texttt{\textbackslash{}paragraph}) should be typeset in
normal case and have a ``.'' at the end.

\paragraph{Internal references.}

We recommend the use of \verb+\autoref+ from \texttt{hyperref}
(automatically loaded by the class).  For instance,
\verb+\autoref{sec:options}+ links to \autoref{sec:options}.

\paragraph{Pictures.}

We recommend the use of the \texttt{tikz} package to render pictures.

In particular, a large variety of crypto pictures made with
\texttt{tikz} is available at \url{http://www.iacr.org/authors/tikz/}.

\paragraph{External pictures.}  The \texttt{graphicx} is loaded by the
class, and is recommended for external figures.

If possible, external figures should be in a vector format: You can
use PDF files when compiling with \texttt{pdflatex}, and EPS files
when compiling with \texttt{latex}, and \texttt{dvips}.  Note that the
\verb+\includegraphics+ command will automatically select a file with
the right extension, so if you write \verb+\includegraphics{figure}+
and have two files \texttt{figure.pdf} and \texttt{figure.eps}, it
should work with both workflows.

\paragraph{Floats.}

All floats should be on the top of the page as shown in \autoref{fig::example}.
%
\begin{figure}
  \center
  \includegraphics[height=2cm]{CC-by.pdf}
  \caption{Example for a floating figure on top of the page.}
  \label{fig::example}
\end{figure}
%
Figure captions should be below the figures, and table captions above
the tables.  The \texttt{float} package loaded by the class should
take care of this automatically.  If want to have several figures side
by side, see the \texttt{[floatrow]} option.
Float captions are complete sentences with a ``.'' at the end.

\paragraph{Tables.}

We recommend the \texttt{booktabs} package to typeset tables as shown in
\autoref{tab::example}.

\begin{table}
  \center
  \begin{tabular}{@{}cc@{}}
    \toprule
    {\bf Column 1} & {\bf Column 2} \\
    \midrule
    Cell 1 & Cell 2 \\
    Cell 3 & Cell 4 \\
    \bottomrule
  \end{tabular}
  \caption{Example for a floating table on top of the page using the
           \texttt{booktabs} package.}
  \label{tab::example}
\end{table}


\paragraph{Algorithms.}

We recommend the \texttt{algorithm}, \texttt{algorithmcx} packages for
algorithms (in particular, \texttt{algpseudocode} for pseudo-code).


\section{For the Editor}

The following commands should be used by the editor to prepare the final
version:

\begin{itemize}
\item \texttt{\textbackslash{}setfirstpage} to set the first page number.
\item \texttt{\textbackslash{}setlastpage} to set the last page number (optional).
\item \texttt{\textbackslash{}setvolume} to set the volume number.
\item \texttt{\textbackslash{}setnumber} to set the edition number.
\item \texttt{\textbackslash{}setDOI} to set the DOI\@.
\item \texttt{\textbackslash{}setPublished} to set the publication date.
\item \texttt{\textbackslash{}setAccepted} to set the notification date.
\item \texttt{\textbackslash{}setRevised} to set the re-submission date for paper that went though major revision.
\item \texttt{\textbackslash{}setReceived} to set the submission date.
\end{itemize}

There is a special \texttt{settings.tosc.tex} file, that sets default values for these commands
and which can be included in the beginning of the main tex file.

\section{Further Information}

More general information can be found in the following documents:
\begin{itemize}
\item General \LaTeX{} documentation, such as the
  \href{http://mirrors.ctan.org/info/lshort/english/lshort.pdf}{(not
    so) short introduction to \LaTeXe};
% \item The \texttt{article} class
%   \href{http://mirrors.ctan.org/macros/latex/doc/clsguide.pdf}{documentation};
\item The \AmS-\LaTeX{}
  \href{http://mirrors.ctan.org/macros/latex/required/amslatex/math/amsldoc.pdf}{documentation}
  and \texttt{amsthm} \href{ftp://ftp.ams.org/pub/tex/doc/amscls/amsthdoc.pdf}{documentation};
\item Documentation of the \LaTeX{} packages used in the class (see below).
\end{itemize}

\subsection{Packages Used}

The class is based on the standard \texttt{article} class, and loads
the following packages:
\begin{itemize}
\item \texttt{geometry}, \texttt{sectsty}, \texttt{fancyhdr},
  \texttt{float}, \texttt{microtype}, \texttt{fontenc}, \texttt{lmodern}
\item \texttt{amsmath}, \texttt{amssymb}, \texttt{amsthm}, \texttt{mathtools}
\item \texttt{graphicx}
\item \texttt{xkeyval}, \texttt{afterpage}
\item \texttt{hyperref}, \texttt{hyperxmp}, \texttt{etoolbox}, \texttt{xcolor} (unless
  the \texttt{[nohyperref]} option is used)
\item \texttt{lineno} (in \texttt{[submission]} mode)
\item \texttt{lastpage} (in \texttt{[final]} mode, unless disabled with \texttt{[nolastpage]})
\item \texttt{floatrow,caption} (with option \texttt{[floatrow]})
\end{itemize}

\section*{Thanks}

We would like to thank people who helped design and improve the class:
Anne Canteaut,
Jérémy Jean,
Marc Joye,
Bart Preneel,
Christian Rechberger,
Tyge Tiessen,
Jonas Wloka.

\section*{Changes}

\begin{description}
\item[v 0.21] First public version
\item[v 0.22] Added documentations.
              Minor tweaks in the class.
\item[v 0.23] More documentation.
              Removed some extra line-numbers with AMS environments in submission mode.
              Make \verb+autoref+ capitalize sections.
              Table caption are now above tables.
              Rewritten running authors and running title.
              Added PDF info (title, author, keyword).
              Optional argument for \verb+\email+.
              Added \texttt{floatrow} option.
\item[v 0.24] Added CC licence text, and added XMP metadata.
              Fixed some metadata transformations.
\item[v 0.26] Added ISSN number and fixed a few bugs with spthm (thanks
  to Marc Joye).
\item[v 0.27] Added paragraph \emph{Helping the editor} in the documentation,
              added \texttt{settings.tosc.tex}.
\item[v 0.90] Moved to github, various fixes.
              Added \texttt{journal} class option to switch between TCHES and ToSC.
              Changed default mode to \texttt{preprint}.
\item[v 0.91] Adding publication in the final version.
              Added new class options: \texttt{notanonymous},
              \texttt{xcolor=xxx}, \texttt{hyperref=xxx},
              \texttt{nolastpage}.
\item[v 0.92] Fixed hyperxmp issues.
              Added template for authors.
\end{description}


\end{document}
